\hypertarget{advanced-functions-callbacks}{%
\chapter{ADVANCED FUNCTIONS \&
CALLBACKS}\label{advanced-functions-callbacks}}

\hypertarget{advanced-functions}{%
\chapter{ADVANCED FUNCTIONS}\label{advanced-functions}}

\hypertarget{variadic-functions}{%
\section{2.1 Variadic Functions}\label{variadic-functions}}

\begin{lstlisting}[language={C++}]
#include <iostream>
#include <cstdarg>
using namespace std;

// Function with variable number of arguments
int sum(int count, ...) {
    va_list args;
    va_start(args, count);
    
    int total = 0;
    for (int i = 0; i < count; i++) {
        total += va_arg(args, int);
    }
    
    va_end(args);
    return total;
}

int main() {
    cout << sum(3, 10, 20, 30) << endl;     // 60
    cout << sum(5, 1, 2, 3, 4, 5) << endl;  // 15
    
    return 0;
}
\end{lstlisting}

\hypertarget{function-recursion}{%
\section{2.2 Function Recursion}\label{function-recursion}}

\begin{lstlisting}[language={C++}]
#include <iostream>
using namespace std;

// Factorial using recursion
int factorial(int n) {
    if (n <= 1) {
        return 1;  // Base case
    }
    return n * factorial(n - 1);  // Recursive case
}

// Fibonacci using recursion (inefficient)
int fibonacci(int n) {
    if (n <= 1) return n;
    return fibonacci(n - 1) + fibonacci(n - 2);
}

// Fibonacci with memoization (efficient)
int fib_memo(int n, int memo[]) {
    if (n <= 1) return n;
    if (memo[n] != -1) return memo[n];
    
    memo[n] = fib_memo(n - 1, memo) + fib_memo(n - 2, memo);
    return memo[n];
}

int main() {
    cout << factorial(5) << endl;  // 120
    cout << fibonacci(10) << endl; // 55
    
    int memo[11];
    for (int i = 0; i < 11; i++) memo[i] = -1;
    cout << fib_memo(10, memo) << endl;  // 55
    
    return 0;
}
\end{lstlisting}

\hypertarget{inline-functions}{%
\section{2.3 Inline Functions}\label{inline-functions}}

\begin{lstlisting}[language={C++}]
#include <iostream>
using namespace std;

// Inline function - compiler may expand code
inline int square(int x) {
    return x * x;
}

// Inline with condition
inline double max_value(double a, double b) {
    return (a > b) ? a : b;
}

int main() {
    cout << square(5) << endl;      // 25
    cout << max_value(3.5, 2.1) << endl;  // 3.5
    
    return 0;
}
\end{lstlisting}

\hypertarget{static-functions}{%
\section{2.4 Static Functions}\label{static-functions}}

\begin{lstlisting}[language={C++}]
#include <iostream>
using namespace std;

// File scope - only visible in this file
static void internal_function() {
    cout << "Internal function" << endl;
}

// Static with counter
int get_call_count() {
    static int count = 0;  // Persists between calls
    return ++count;
}

int main() {
    cout << get_call_count() << endl;  // 1
    cout << get_call_count() << endl;  // 2
    cout << get_call_count() << endl;  // 3
    
    internal_function();  // OK in same file
    
    return 0;
}
\end{lstlisting}

\begin{center}\rule{0.5\linewidth}{0.5pt}\end{center}

\hypertarget{function-pointers-callbacks}{%
\chapter{FUNCTION POINTERS \&
CALLBACKS}\label{function-pointers-callbacks}}

\hypertarget{function-pointers}{%
\section{3.1 Function Pointers}\label{function-pointers}}

\begin{lstlisting}[language={C++}]
#include <iostream>
using namespace std;

// Function pointer declaration: return_type (*name)(parameters)
int add(int a, int b) {
    return a + b;
}

int subtract(int a, int b) {
    return a - b;
}

int multiply(int a, int b) {
    return a * b;
}

int main() {
    // Declare function pointer
    int (*operation)(int, int);
    
    // Assign function to pointer
    operation = add;
    cout << operation(5, 3) << endl;  // 8
    
    operation = subtract;
    cout << operation(5, 3) << endl;  // 2
    
    operation = multiply;
    cout << operation(5, 3) << endl;  // 15
    
    return 0;
}
\end{lstlisting}

\hypertarget{function-pointer-arrays}{%
\section{3.2 Function Pointer Arrays}\label{function-pointer-arrays}}

\begin{lstlisting}[language={C++}]
#include <iostream>
using namespace std;

int add(int a, int b) { return a + b; }
int subtract(int a, int b) { return a - b; }
int multiply(int a, int b) { return a * b; }
int divide(int a, int b) { return a / b; }

int main() {
    // Array of function pointers
    int (*operations[4])(int, int) = {
        add, subtract, multiply, divide
    };
    
    int a = 20, b = 4;
    
    for (int i = 0; i < 4; i++) {
        cout << "Result: " << operations[i](a, b) << endl;
    }
    // Output: 24, 16, 80, 5
    
    return 0;
}
\end{lstlisting}

\hypertarget{callbacks}{%
\section{3.3 Callbacks}\label{callbacks}}

\begin{lstlisting}[language={C++}]
#include <iostream>
#include <vector>
using namespace std;

// Callback function type
typedef void (*Callback)(const string&);

class Button {
private:
    Callback on_click;
    
public:
    void set_click_handler(Callback callback) {
        on_click = callback;
    }
    
    void click() {
        if (on_click) {
            on_click("Button clicked!");
        }
    }
};

void handle_click(const string& message) {
    cout << message << endl;
}

int main() {
    Button button;
    button.set_click_handler(handle_click);
    button.click();  // Output: Button clicked!
    
    return 0;
}
\end{lstlisting}

\begin{center}\rule{0.5\linewidth}{0.5pt}\end{center}
